% Part 3 compression at depth (the director, 2026-11-20; DECISION-20261120a item 5), director-helper-3: sentences moved
% verbatim from the note's Part 3, in order, one paragraph per source paragraph (move manifest H3-P3C-*). References
% into the note resolve through main_note.aux (xr).
\section{Part 3 detail moved from the note}
This section holds the construction detail, qualifier detail and secondary checks that the compressed Part~3 no longer prints, as the note printed them.

\subsection{From the uncertainty of the two-dimensional cross section (\S\ref{sec:systematics})}
% From sec_systematics.tex, paragraph 1
The universe, statistical and training-seed blocks are computed with the LightGBM estimator of the quoted 2D central value (\S\ref{sec:method-config}); the universe and statistical blocks use its training seed (seed~42), and the training-seed block uses five seed settings that include it. A 3D budget built with the analogous 3D grid construction is not quoted, and no 4D uncertainty magnitude is reported (\S\ref{sec:4d}). App.~\ref{app:statmethods} gives the full mathematics and worked examples using the reported numbers (N-D declarations in App.~\ref{app:statmethods-nd}). The resulting budget and its agreement with the published total are in \S\ref{sec:results-uq}; the conventions the covariance follows and the model-dependence component reported beside it are in \S\ref{sec:uq-2d-scope}, and what it omits and its validation status are in \S\ref{sec:limitations-2d}.

% From sec_systematics.tex, paragraph 2
Each universe is re-unfolded with LightGBM at seed~42 and compared with the quoted 2D central value, unfolded with the same estimator and seed. The five lateral bands, which shift reconstructed kinematics rather than event weights, are taken from re-unfolds of re-selected universes (\S\ref{sec:uq-higherdim}) at the same seed. The band sum omits the nonresonant-pion band Rvp1pi, which encodes the same shift as Rvn1pi (\S\ref{sec:uq-2d-scope}); it holds 43 bands plus the flat normalization. App.~\ref{app:pair} gives the per-band covariance.

% From sec_systematics.tex, paragraph 3
Figure~\ref{fig:uqbands} shows the fractional uncertainty grouped as in Figs.~6--7 of the published measurement; the full $205\times205$ covariance is projected exactly, $C_{1D}=PC_{2D}P^{\top}$, before division by the 1D central value.

\paragraph{Flux $1/\Phi$ normalization.}
% From sec_systematics.tex, paragraph 4
Varying $\Phi$ per flux universe in the dominant $\sigma\propto1/\Phi$ term gives a Flux band of \SI{\fluxBand}{\percent} with the constraint weights, and a combined budget, from our covariance alone, of \SI{\uqMedian}{\percent}. Target-nucleon normalization contributes a fully-correlated rank-1 band of \SI{1.4}{\percent}. The unweighted 100-universe flux covariance has participation-ratio effective rank $\approx1$. Its leading mode correlates at $\rho=\rhoFluxMuon$ with the rank-1 muon-energy-scale band of the joint flux and muon-energy-scale block of Ref.~\cite{Bashyal:2021tqu}, as rederived here. The quoted budget does not include this cross block: with the MINOS universes divided by their own flux, the correlated-flux effect is already in the MINOS band, and adding the block would count it twice.

% From sec_systematics.tex, paragraph 5
Each replica draws its weights in three independent streams, for data, signal simulation and background simulation; possible shared events between the two simulations were not assessed. Because the background simulation is resampled in every replica, the separable background-only term is contained in the statistical block. The reconstructed-level and truth-level weights of a simulated event share one draw, which preserves within-event migration correlations. Fixing the GBDT seed per replica separates bootstrap and training-seed variances. The replica mean differs from the central value by a per-bin median of 4.37 standard errors of the mean ($\sigma/\sqrt{300}$), or 0.25 of the statistical standard deviation. Because the block is the covariance about the replica mean, it does not include this offset.
% absence-ok: "not assessed" is the block's own construction scope, unchanged by RE-2; the sentence edited here is W4 (copied with the moved sentence, H3-P3C)

% From sec_systematics.tex, paragraph 7
The quoted central value is one unfold at one seed. Its statistical block is the bootstrap at that seed and already contains the training noise at that seed, so the added component is this covariance of the per-seed means. A spread of single unfolds on fixed data is a different statistic. With four degrees of freedom, its variance has a Gaussian relative standard error of about \SI{71}{\percent}. The bootstrap\,$\times$\,seed interaction and the universe\,$\times$\,retraining term are reported and not added (\S\ref{sec:limitations-2d}).

% From sec_systematics.tex, paragraph 9
For two bands whose universes barely change the unfolding inputs, EtaNCEL and NormNCRES, the responses are consistent with instability of the GBDT training contributing at the scale of the training-seed variation (\S\ref{sec:uq-ml}). They neither establish the mechanism for every small band nor justify subtracting an overlap from the budget.

% From sec_systematics.tex, paragraph 10
The checks of the block's construction are in App.~\ref{app:devhistory-2dchecks}.

% From sec_systematics.tex, paragraph 12
After unfolding, the double ratio (data over simulated truth) of the data-taking periods (playlists) 1G--1M to the earliest periods, 1A--1F, has a $\pPar$ step of $(+6.89\pm0.46)\%$ at \SI{8}{GeV/c}. The step is already present in the reconstructed data, at $+\fvval{rp-pb-reco}\,\%$ relative to 1A--1F (data minus background over each period's signal simulation, scaled to its exposure; Fig.~\ref{fig:runperiod-steps}(a)). It therefore survives unfolding at about its reconstruction-level size. The 1N--1P group shows a smaller step, $+\fvval{rp-pc-reco}\,\%$ reconstructed and $+(\fvval{rp-pc-unf}\pm\fvval{rp-pc-unf-se})\,\%$ unfolded, and lies above 1A--1F in every $\pPar$ bin below \SI{8}{GeV/c}, in central value, by \SIrange{\fvval{rp-pc-lev-reco-lo}}{\fvval{rp-pc-lev-reco-hi}}{\percent} at reconstruction level and \SIrange{\fvval{rp-pc-lev-unf-lo}}{\fvval{rp-pc-lev-unf-hi}}{\percent} after unfolding, an excess of unknown cause. Diluted into the combined result, the raw shift of those playlists is small in each bin: every marginal $z$ is below 0.7. The shift's $\pPar$ $\chi^{2}$ against the earlier 2D budget is nevertheless 34.4, where noise would give 2.0.

% From sec_systematics.tex, paragraph 13
The block is built from separate seed-42 unfolds of three run periods, the playlist groups 1A--1F, 1G--1M and 1N--1P, each unfolded through its own simulated detector response. The block uses each period's unfolded result over that period's own simulated truth. Each period is compared with the exposure-weighted mean of the three. This avoids choosing a reference period. The spread, applied to the central value, gives the block. It is added after the other blocks. Its per-bin median fractional standard deviation is \SI{2.09}{\percent} (84th percentile \SI{3.34}{\percent}, maximum \SI{6.48}{\percent}). The fractional standard deviation is \SIrange{1.17}{3.36}{\percent} on the $\pPar$ projections and \SI{1.97}{\percent} on the integrated cross section. Table~\ref{tab:pl-added} gives its effect on the total. Taking 1A--1F as the reference would give a smaller block (per-bin median \SI{1.38}{\percent}, integral \SI{1.09}{\percent}).

% From sec_systematics.tex, paragraph 14
Because the block is built from the observed differences among the periods, it contains them by construction (App.~\ref{app:devhistory}). That containment gives no information about coverage or whether the size is right; the block's other limits are stated in \S\ref{sec:limitations-2d}. For the same reason, the double-ratio $\chi^{2}$ against any total that contains the block is bounded whatever the data. It does not test the block. The raw 1G--1M $\pPar$ $\chi^{2}$ is 33.8 against the other blocks alone and 2.60 with the block added. These values indicate sensitivity only. % G2R: RE-2-shapes run 2 (C_new), g2r-integration/ledger.tsv g2r-cpl-raw-vs-s5, g2r-cpl-raw-vs-q; a source comment citing RE-2 part C for this value gives the earlier value

\paragraph{Status of the 2D result.}
% From sec_systematics.tex, paragraph 15
Section~\ref{sec:limitations-2d} states its validation status and what the covariance does not include. That section also gives the seed-related terms reported but not added, the inputs carried without confirmation and the uncorrected bias from pooling the playlists. The dependence on the generator shape of the prior is reported beside the covariance as $C_{\mathrm{model}}$ (below).

\paragraph{Generator-shape model dependence.}
% From sec_systematics.tex, paragraph 16
For each of three generators $k$ (GENIE with meson-exchange currents, NuWro and GiBUU), the prior's truth distribution is reweighted by the ratio of that generator's $(\pT,\pPar)$ shape to GENIE's, with its normalization kept. The ratio is constant on each reported $(\pT,\pPar)$ cell. A relative-error gate, fixed before the reweighted unfolds were run, decides each cell's ratio. A cell whose own ratio passes the gate keeps it. Otherwise the cell takes the ratio pooled over a coarser cell, fixed before the measurement, when that pooled error passes the same gate. When neither passes, the cell leaves the prior unchanged. A row outside the grid, an unreported cell, and GiBUU where that sample is not defined also leave the prior unchanged. The data are unfolded again with the central value's estimator; $s_k$ is the per-cell relative change of the result. That sum is a fractional covariance. In cross-section units it is $\mathrm{diag}(x)\,C_{\mathrm{model}}\,\mathrm{diag}(x)$, with $x$ the cell-integrated cross sections of this result.

% From sec_systematics.tex, paragraph 17
Each generator enters as one fully correlated component, unscaled, so $C_{\mathrm{model}}$ has rank three. On the integrated cross section $C_{\mathrm{model}}$ gives \SI{\cmodelIntegral}{\percent}, and the total stays at \SI{\uqIntegral}{\percent}: the reweightings change only the prior's shape, so $C_{\mathrm{model}}$ does not test the model dependence of the normalization. The comparison with the published uncertainty (\S\ref{sec:results-uq}) uses the combined covariance alone, because the published uncertainty has no component for this dependence either. % G2R: RE-2-shapes run 2 (C_new), g2r-integration/ledger.tsv g2r-qr-cmodel-p84, g2r-qr-cmodel-max; a source comment citing RE-2 part C for this value gives the earlier value

% From sec_systematics.tex, paragraph 18
Where a pooled cell, a noise-gated cell, or an unreported cell meets another cell, the construction can draw a step of more than 0.1 in a generator's reweighting ratio; the cells beside such steps enter with their measured shifts. Identified after the measurement, these 52 cells hold \SI{59.4}{\percent} of the trace of $C_{\mathrm{model}}$ and carry most of its large shifts (App.~\ref{app:devhistory}).

\paragraph{Conventions.}
% From sec_systematics.tex, paragraph 19
This follows the MAT treatment of universes (\S\ref{sec:uq-universes}), which we use as a stated convention. Across bands, the per-bin median offset of a band's center from the central value is \SIrange{0.12}{0.57}{\percent} of the central value, or 0.09--2.85 times the band's $\sigma$. The statistical block is the covariance about the replica mean (\S\ref{sec:uq-boot}). % G2R: RE-2-shapes run 2 (C_new), g2r-integration/ledger.tsv g2r-centring-min-rel, g2r-centring-max-rel; a source comment citing RE-2 part C for this value gives the earlier value

\paragraph{Missed events in the universe unfolds.}
% From sec_systematics.tex, paragraph 20
In the systematic-universe unfolds, each event that fails reconstruction carries that universe's own truth weight. An earlier construction, since replaced, is described in App.~\ref{app:devhistory-2dchecks}. % D1 follow-on (the director's go of 2026-11-16): DECISION-20261114-joseph-d1-g2r-adopted.md items 2-3; COV2D-G2R-DECLARATION.md sec. 10, VALID row (\"removed. The covariance text states that the systematic-universe unfoldings weight each missed event at its universe's truth weight\"). The earlier approximation's two-universe tests are in the ledger: history-cuts/sec_systematics.md S7.

\paragraph{Comparison with the earlier 2D budget.}
% From sec_systematics.tex, paragraph 21
The quoted budget corrects three defects in the construction of the earlier 2D budget: a simulated-truth input of its systematic sweep that differs from that of its reference, a nonresonant-pion shift counted twice, and background-template statistics that were not propagated. None of the three is present in the quoted budget. The records document gives each defect with its diagnostics.

\subsection{From the two-dimensional cross section (\S\ref{sec:results})}
% From sec_results.tex, paragraph 3
The comparison leaves out the generator-shape component $C_{\mathrm{model}}$ (\S\ref{sec:uq-2d-scope}), because the published uncertainty has no such component either. With $C_{\mathrm{model}}$, the per-bin median is \SI{\qrCmodelMedian}{\percent}. The systematic-universe block accounts for a median 0.87 of the per-bin variance and 0.83 of the integral's variance, and the run-period block (\S\ref{sec:uq-combine}) raises the per-bin median from \SI{\uqSFiveMedian}{\percent} to \SI{\uqMedian}{\percent}.

% From sec_results.tex, paragraph 4
Our flux band carries the per-universe $\nu$--$e$ and inverse-muon-decay constraint weights of the MINERvA Analysis Toolkit. The statistical and ML blocks together contribute under a tenth of the combined $\sqrt{\mathrm{Tr}\,C}$. Every block is computed with the LightGBM estimator of the quoted central value (\S\ref{sec:method-config}). Relative uncertainties are taken with respect to that central value, which is the universe sweep's central value at training seed~42.

% From sec_results.tex, paragraph 7
With the covariance of this result alone, $\chi^{2}$ is computed only in the covariance's well-constrained subspace, with a truncated spectral pseudo-inverse. The per-bin pulls and the truncated result are therefore quoted together, and an unconstrained pseudo-inverse is not used as a goodness-of-fit test on its own. The regularization scan (App.~\ref{sec:rank}) was made for the exact-split cross-check on an earlier covariance and was not repeated for the quoted one.

% From sec_results.tex, paragraph 8
Restricted to the 73 best-measured directions, $\chi^{2}=136$ ($\chi^{2}/205=0.66$; the divisor is the number of reported bins, not the number of directions kept). No single $\pT$ or $\pPar$ band carries more than \SI{15.2}{\percent} of the total.

% From sec_results.tex, paragraph 12
A value near one does not by itself indicate reproduction. What the benchmarks leave open is stated in \S\ref{sec:limitations-2d}.

% From sec_results.tex, paragraph 13
It is farther by \lrval{pp-kbar-delta} per direction, a fraction \lrval{pp-kbar-frac} of the three-group distance, beyond noise estimated from \lrval{pp-nrep} replicas.

% From sec_results.tex, paragraph 14
The two-dimensional estimator with deterministic training, the five-observable estimator's training mode, agrees with the two-dimensional central value within its paired noise, narrowly: $T=\fvval{c1-T}$ over \fvval{c1-nrep} paired replicas, \SI{\fvval{c1-below-pct}}{\percent} below its 95th-percentile threshold of \fvval{c1-q95}.

% From sec_results.tex, paragraph 15
Our result tracks the data far more closely than the model does (Fig.~\ref{fig:modelmaps}). The published analysis~\cite{MINERvA:2021owq} also compares with NuWro~\cite{Golan:2012wx}, GiBUU~\cite{Buss:2011mx} and NEUT~\cite{Hayato:2009zz}. OmniFold's per-event reweighted truth ensemble can be compared with any such prediction in \emph{any} projection of the classifier features without re-unfolding. An observable outside that space requires a new unfolding, as in the 3D extension (\S\ref{sec:3d}), which adds $\Eavail$ as a classifier feature.

\subsection{From the validation of the two-dimensional result (\S\ref{sec:validation})}
% From sec_validation.tex, paragraph 1
In a simulated experiment, \emph{truth} is the known target distribution, the \emph{prior} is the starting simulated truth distribution, and the \emph{response} describes reconstruction and selection conditional on truth. Bias is the mean estimate minus that target over repeated experiments. A single closure residual also contains sampling fluctuations. Repeating a numerical calculation checks neither the coverage and width of the intervals nor the adequacy of the physical model.

\paragraph{Closure.}
% From sec_validation.tex, paragraph 2
The alternative-model test trains on the central-value simulation and takes its pseudo-data from an alternative systematic universe.

\paragraph{Injected-feature recovery.}
% From sec_validation.tex, paragraph 3
The test requires the residual left after recovering a deviation from the prior to be far smaller than the deviation itself. Unfolding cannot create discriminating power (3D: \S\ref{sec:3d-closure}).

% From sec_validation.tex, paragraph 4
A data-versus-prior $\chi^2$ with a diagonal covariance of statistical variances only is larger at truth level than at reco level. This is expected at high resolution: detector smearing lowers the reco-level value, and the diagonal covariance omits the unfolding covariance. The comparisons that use a full covariance are the $\chi^2$ against the published result (\S\ref{sec:results-tension}) and the generator comparisons.

% From sec_validation.tex, paragraph 6
The hidden-variable closure does change the effective response: it reweights reconstructed events by a function of $\Delta\pT$ and leaves the truth weights unchanged (\S\ref{sec:closure-hv}). Appendix~\ref{app:response-mismatch} defines a separate test, a closure that rescales only the reconstructed $\Eavail$ of the pseudo-data, together with a more complete leading-nucleon follow-up. No result of that reconstruction-only test is reported here; the same appendix reports noise-free closures with a shifted recoil response on the sixteen cells, which have been run. The five-dimensional estimator, including its background subtraction, is evaluated under the technical closure checks of \S\ref{sec:val-fiveaxis}. Its coverage is unvalidated.

% From sec_validation.tex, paragraph 10
Each toy draws pseudo-data event by event as Poisson counts with means equal to the simulated reconstructed weights. It resamples the simulation with the per-event Poisson(1) bootstrap of the statistical band and unfolds with that band's estimator, LightGBM with five iterations at a fixed training seed, the learner of the quoted central value at a different seed (\S\ref{sec:method-config}). The reference $T$ is the unfluctuated simulated truth, identical in every toy. The no-fluctuation configuration reproduces both the production closure output and $T$ exactly.

% From sec_validation.tex, paragraph 11
In reported bin $b$ of toy $t$, $z_{tb}=(U_{tb}-T_b)/(r_b T_b)$, where $r_b$ is the production band's relative spread. Coverage $C_k$ is the fraction of all (toy, bin) pairs with $|z|\le k$, with intervals from resampling whole toys. The band counts as nominal if its ratio to the true scatter lies in $[0.9,1.1]$, which gives the windows $C_1\in[\covTwoDWinOneLo,\covTwoDWinOneHi]$ and $C_2\in[\covTwoDWinTwoLo,\covTwoDWinTwoHi]$. A failure at the interim analysis of 100 toys, with 99.5\% intervals, stops the study.

% From sec_validation.tex, paragraph 12
The pull RMS is $\covTwoDPullRMS$, and the mean pull is consistent with zero.

% From sec_validation.tex, paragraph 13
The positive control rescales the band on the same toys. At $\times0.7$, $C_2$ moves to $\covTwoDCTwoLowCtl$, below its window. At $\times1.3$, $C_2$ moves to $\covTwoDCTwoHighCtl$, inside its window, and $C_1$ to $0.785$, above its window. The test flags both directions.

% From sec_validation.tex, paragraph 14
App.~\ref{app:devhistory} gives that analysis, its attribution rule and the rebuild.

% From sec_validation.tex, paragraph 17
A pseudo-experiment is inside a cell when the difference between its unfolded result and the held-out truth lies within that cell's standard deviation under the remade covariance. It is inside the joint region when its $\chi^{2}$ on that covariance, with the degrees of freedom printed below, is at most the $\chi^{2}$ quantile at the one-standard-deviation Gaussian content. This containment is why the test is called a coverage test. The mean $\chi^{2}/\mathrm{ndf}$ of the pseudo-experiments is $\lrval{cov2d-ar23-chi2ndf}\pm\lrval{cov2d-ar23-chi2ndf-se}$ ($\mathrm{ndf}=\lrval{cov2d-ar23-ndf}$) at \fvval{f1-tune-ar23} and $\lrval{cov2d-g18-chi2ndf}\pm\lrval{cov2d-g18-chi2ndf-se}$ ($\mathrm{ndf}=\lrval{cov2d-g18-ndf}$) at \fvval{f1-tune-g18}. The test's false-flag rate under its null's law is at most \lrval{cov2d-ar23-size} and \lrval{cov2d-g18-size} per tune. In a descriptive study made after the test's result, truncating the joint test to its \lrval{t2a-r1} or \lrval{t2a-r2} largest-eigenvalue directions still finds none of the \lrval{cov2d-ntoys} pseudo-experiments in its region at either tune, so truncation to the better-determined directions does not remove the joint under-coverage. In the same study, the mean $\chi^{2}$ splits into a part from the pseudo-experiments' mean error $\bar e$, $\bar e^{\top}C^{+}\bar e$, and a part from their spread $S$, $\mathrm{tr}(C^{+}S)$, with $C$ the remade covariance. At \fvval{f1-tune-ar23} and \fvval{f1-tune-g18}, the mean-error part is \fvval{t2a-bias-share-ar23}\,\% and \fvval{t2a-bias-share-g18}\,\% of the mean $\chi^{2}$, or \fvval{t2a-bias-ndf-ar23} and \fvval{t2a-bias-ndf-g18} per degree of freedom. The spread part is \fvval{t2a-spread-ndf-ar23} and \fvval{t2a-spread-ndf-g18} per degree of freedom, where 1 would mean that the covariance describes the spread. In the covariance's eigenbasis, the \fvval{t2a-ntop} directions with the largest eigenvalues carry \fvval{t2a-top-ar23}\,\% and \fvval{t2a-top-g18}\,\% of the mean-error part, and the directions beyond the \fvval{t2a-r1} largest carry \fvval{t2a-beyond1-ar23}\,\% and \fvval{t2a-beyond1-g18}\,\%. The decomposition locates the $\chi^{2}$ and leaves its cause open. These numbers are the test's regrade against the adopted 2D covariance (\S\ref{sec:uq-2d-scope}); the regrade leaves the verdict, the joint counts and their thresholds, and the flagged cells and their in-interval counts as the first grade had them (App.~\ref{app:devhistory}).

% From sec_validation.tex, paragraph 19
It adds to the fixed-truth result of \S\ref{sec:coverage}; it does not replace it.

\paragraph{Unbinned diagnostic (C2ST).}
% From sec_validation.tex, paragraph 20
On this training sample, that drop is a discrimination diagnostic. It does not show that the mismatch was removed, and it is not evidence of generalization or of calibrated agreement. As run, the C2ST trains and evaluates on the same sample without cross-fitting. A valid null would require re-running the full cross-fitted OmniFold pipeline per pseudoexperiment. Calibrated unbinned goodness-of-fit remains open~\cite{Canelli:2025guide}.
